\documentclass[10pt]{article}
\usepackage[letterpaper,margin=0.68in]{geometry}
\usepackage{amsmath,amssymb}
\usepackage{enumitem}
\usepackage{xcolor}
\usepackage{fancyhdr}
\usepackage{microtype}
\usepackage[most]{tcolorbox}
\definecolor{olive}{HTML}{4D5430}
\definecolor{gold}{HTML}{9A7B22}
\definecolor{pale}{HTML}{F5F1DC}
\definecolor{softgreen}{HTML}{EDF0E2}
\pagestyle{fancy}\fancyhf{}\lhead{\textcolor{olive}{MATH\& 107 Answer Key}}\rhead{\textcolor{gold}{Fall 2026}}\cfoot{\thepage}
\setlength{\headheight}{14pt}
\setlength{\parindent}{0pt}\setlength{\parskip}{5pt}
\setlist[enumerate]{leftmargin=*,itemsep=5pt,topsep=4pt}
\newtcolorbox{solbox}[1]{colback=pale,colframe=olive,boxrule=.7pt,arc=2mm,title=\textbf{#1},fonttitle=\color{olive},left=2mm,right=2mm,top=1.5mm,bottom=1.5mm}
\newcommand{\modulekey}[2]{\clearpage{\Large\bfseries\textcolor{olive}{Module #1: #2 — Answer Key}}\par\textcolor{gold}{\rule{\linewidth}{1.2pt}}\par}
\begin{document}
\modulekey{1}{Number Systems}
\textbf{Prequiz}
\begin{enumerate}
\item \(-35\) is an \textbf{integer}. It is not natural, and integers are contained in both the rationals and reals, so \(\mathbb Z\) is the smallest standard set listed that contains it.
\item The trip is
\[
-6+10+(-14)=-10.
\]
The final position is \(\boxed{-10}\).
\item Cross-products:
\[
6\cdot21=126,\qquad 9\cdot14=126,
\]
so the fractions are equivalent. Also,
\[
\frac69=\frac23,\qquad \frac{14}{21}=\frac23.
\]
\item A suitable recursion is
\[
H_0=800,\qquad H_{n+1}=H_n+120.
\]
After four years,
\[
H_4=800+4(120)=\boxed{1280}.
\]
\item
\[
B_0=800,\qquad B_{n+1}=1.10B_n.
\]
Then
\[
B_3=800(1.10)^3=\boxed{1064.8}.
\]
The \textbf{ratio} is constant: each year's value is \(1.10\) times the previous year's value.
\end{enumerate}

\textbf{Matching quiz}
\begin{enumerate}
\item \(-8\) is an \textbf{integer}.
\item
\[
-3+11+(-6)=\boxed{2}.
\]
\item
\[
\frac8{12}=\frac23,\qquad \frac{18}{27}=\frac23,
\]
so the fractions are equivalent.
\item
\[
H_0=1200,\qquad H_{n+1}=H_n+75,
\]
so
\[
H_5=1200+5(75)=\boxed{1575}.
\]
\item
\[
B_0=1200,\qquad B_{n+1}=1.05B_n,
\]
so
\[
B_2=1200(1.05)^2=\boxed{1323}.
\]
\end{enumerate}
\end{document}
